\section{[Remove?] Virtualization of $\FF_2$-linear combinations lifted to integers}\label{s:virtualization}

\albertimportant{Do we really need this section? Isn't the technical overview enough? Having contents only in the TO is a bit weird academically speaking, but this section is really a straightforward rewriting of the corresponding TO section, with some more formal symbols and language}


In this section we expand the core relation $\LinBitsDual$ to impose, instead of constraints of the form $f=\Enc_\cC\circ \pack\circ \pi_2(\bff)$ and $\langle \pievalfield(\bff), \vv  \rangle = \mu$, where $\bff\in \bits^{\codedim} \subseteq \ZZ^{\codedim}$, constraints as $\Enc_\cC\circ\pack\circ\pi_2(\bff) = f$  and $$\langle \pievalfield\circ \pi_2^{-1}(\Mmatrix\cdot \pi_2(\bff)) \rangle = \mu$$ for a matrix $\Mmatrix\in \FF_{2}^{\codedim' \times \codedim}$. Precisely, given a new parameter $\codedim'$, we define 
%
\begin{align*}
&\LinBitsDual^\virtual(\bits,\pi_q, \evalfield,\cC)\\ &\qquad \qquad \qquad \defeq \left\{ \left(\begin{aligned} &\inp=(\mu, \vv, \Mmatrix)\in \evalfield \times \evalfield^{\codedim'} \times \FF_2^{\codedim' \times \codedim}, \\&\oracle = f \in \comfield^\codelength, \\ &\wit= \bff\in \bits^\codedim \subseteq \ZZ^\codedim\end{aligned}\right) \;\middle|\;  \begin{aligned} &f=\Enc_{\cC}\circ \pack \circ \picomfield(\bff)\\ \wedge \  & \langle \pi_q\circ \pi_2^{-1}(\Mmatrix \cdot \pi_2(\bff)), \vv\rangle  = \mu  \end{aligned} \right\}
\end{align*}
\albert{We could replace $M\cdot \pi_2(\bff)$ by the output of a GKR circuit on $\pi_2(\bff)$}
%
The relation has the tag ``virtual'' to reference the intuition that  $\LinBitsDual^\virtual(\bits,\pi_q, \evalfield,\cC)$ expresses a linear constraint on the vector $\bfh=\pi_q\circ \pi_2^{-1}(\Mmatrix \cdot \pi_2(\bff))\in \evalfield^{\codedim'}$, which is not committed directly: only $\pi_2(\bff)$  is. Note that $\bfh$ is a linear expression over $\comfield$ of the elements in $\pi_2(\bff)$, lifted to the integers, and then projected to $\evalfield$.  %In \cref{?} we arithmetize SHA-256 as an instance of $\LinBitsDual^\virtual(\bits,\mathsf{id}, \ZZ,\cC)$ 



\begin{construction}[IOP for $\LinBitsDual^\virtual$]\label{c:virtual_iop}
\textbf{Parameters.} $(\evalfield,\ringmorphism,\comfield,\codedim',\codedim_1,\codedim_2, \codedim, \cC)$ comprised of: a ring $\evalfield$ (either $\ZZ$ or a finite field); a ring homomorphism $\ringmorphism:\ZZ \to \evalfield$; a binary field $\comfield$; a virtual witness length $\codedim'$; tensor decomposition dimensions $\codedim_1, \codedim_2$ satisfying $\codedim_1\cdot \codedim_2 = \codedim'$; a non-virtual witness length $\codedim$; a linear code $\cC$ over $\comfield$ of dimension $\codedim/\comfieldexp$ and blocklength $\codelength$.\\ 

\noindent \textbf{Inputs.}
The prover and verifier receive $\inp=(\mu, \vv,\Mmatrix)\in \evalfield \times \evalfield^{\codedim'} \times \FF_2^{\codedim'\times\codedim}$ and  $\oracle= f \in \comfield^\codelength$. The prover receives $\wit= \bff\in \bits^\codedim \subseteq \ZZ^\codedim$ such that $(\inp,\oracle, \wit) \in \LinBitsDual^\virtual(\bits,\pi_q, \evalfield,\cC)$.  %We assume $M$ is of the form described above.



 \begin{enumerate}[leftmargin=*]
     \item Define  $\bfh = \pi_2^{-1}\left(\Mmatrix\cdot \pi_2(\bff)\right) \in \{0,1\}^{\codedim'} \subseteq \ZZ^{\codedim'}$. %By \eqref{e:M_pi_f}, $$\bfh=(\bfh_{*j})_{j\in [\codedim_2]} = \left(\pi_2^{-1}\left( M_0\cdot \pi_2(\bff_{*j}) \right)\right)_{j\in [\codedim_2]},$$ with $ \bfh_{*j}\in \bits^{\codedim_1'}$. Let $f= \Enc_{\cC}\circ \pack\circ\pi_2(\bff)\in \bits^{\codelength} \subseteq \ZZ^{\codelength}$. 
 %
 Prover and verifier execute \cref{c:core_iop} until before Phase 3, with parameters\footnote{\cref{c:core_iop} would require a code $\cC$ as parameter whose dimension matches $\codedim'$. This is not the case in our current setting, but since \cref{c:core_iop} is run only up to before Phase 3 (and so the code $\cC$ is not used in this partial run), this is not an issue.}
     $$
     (\evalfield,\comfield, \codedim', \codedim_1, \codedim_2, \cC)
     $$ 
     and inputs 
    $$
    \inp' = (\mu, \vv),  \quad \oracle'=f,\quad  \wit' = \bfh \in \{0,1\}^{\codedim'} \subseteq \ZZ^{\codedim'}.
    $$
     Note that until that point, the verifier has not made any queries to $\oracle'$. So, in terms of completeness, it does not matter, yet, that $f$ is not the encoding of $\pi_2(\bfh)=M\cdot \pi_2(\bff)$.
 \item  At the end of Phase 3 of \cref{c:core_iop}, prover and verifier are left with the claims 
\begin{align*}\label{e:iopp_mle_claims_2}
&\left((\vv, \mu'), \ f, \ \picomfield(\bfh)\right) \in \LIN( \comfield,\cC,\pack) \\& \qquad\qquad\qquad = \left\{ \left(\begin{aligned} &\inp=(\uu,\mu)\in \comfield^\inplen, \\&\oracle = g \in \comfield^{\codelength}, \\ &\wit= \bfg\in (\comfield^\codedim)\end{aligned}\right) \;\middle|\;  \begin{aligned} 
&g=\Enc_{\cC}\circ \pack (\bfg) \\    &\langle \uu, \bfg\rangle = \mu  \end{aligned}\right\}
\end{align*} for some vector $\vv\in\comfield^{\codedim'}$ \albert{$\vv$ was already used in the input} and field element $\mu'\in \comfield$ output during the IOR $\Pi_{\mathsf{red}}$. 

As we mentioned, this claim as is is false, because $f\neq \Enc_{\cC}\circ \pack\circ \pi_2(\bfh)$ --indeed, the dimension of $\cC$ does not match that of $\pack\circ \pi_2(\bfh)$. Instead of, as done in \cref{c:core_iop}, using an IOP for $\LIN( \comfield,\cC,\pack)$ to attempt to prove the above claim, prover and verifier proceed as follows.

Observe that, if the prover is honest,
%
$$
\mu'=\langle \vv, \pi_2(\bfh)\rangle = \langle \vv, \Mmatrix\cdot \pi_2(\bff) \rangle = \langle \Mmatrix^{\mathsf{T}}\cdot \vv, \pi_2(\bff)\rangle,
$$
%
where $\mathsf{T}$ denotes transposition. Hence the following claim is true:
%
\begin{equation}\label{e:virtualization_final_claim}
((\Mmatrix^{\mathsf{T}}\cdot \vv, \mu'), f, \pi_2(\bff)) \in \LIN(\comfield, \cC, \pack).
\end{equation}
%
Now prover and verifier use an IOP for $\LIN(\comfield, \cC, \pack)$ to prove \eqref{e:virtualization_final_claim} \albert{comments on verifier?}. 
\end{enumerate}
%
\end{construction}

\begin{remark}
\albert{Comments on how the verifier handles $\Mmatrix^{\mathsf{T}}\cdot \vv$}
\end{remark}

\begin{theorem}[Completeness and round-by-round knowledge soundness of \cref{c:virtual_iop}]\label{t:virtual_iop} 
\albertimportant{Should we be more formal with this theorem?} \albert{Completeness missing}
   Under the same hypothesis as \cref{t:thm_core_IOPP} for $\delta,\basegenerator, \codedim',\codedim_1, \codedim_2, \Pi_\red, \Pi_\iopp$,  \cref{c:virtual_iop} is a perfectly complete IOP for $\LinBitsDual^{\virtual}$  and has round-by-round knowledge soundness with relaxation $((\LinBitsDual^{\virtual})^{\delta}, \bot)$, with errors $$\kappa_0^*, (\ksred_i)_{i\in [t_\red]}, (\ksiopp_i)_{i\in [t_\iopp]}$$ as defined in \cref{t:thm_core_IOPP}.  
   The errors $\kappa_{\red,i}$ correspond to the errors of  the IOR $\Pi_\red$ when used on triples of the form $$(\inp_\red=((g^{\bm{\gamma}_i})_{i\in [\codedim_1]}, (g^{\mu_j})_{j\in [\codedim_2]}), \oracle_\red=f, \wit_\red =\pi_2(\bfh))\in \PESAT(\comfield,Q_{\mathbb{B}\text{it}\ZZ}, \pack)$$ for certain $\bm{\gamma}_i$ and $\mu_j$ (so, the witness has dimension $\codedim'$ and we use the decomposition $\codedim' = \codedim_1 \cdot \codedim_2$). The errors $\kappa_{\iopp,i}$ correspond to the errors of  the IOP $\Pi_\iopp$ when used on triples of the form $$(\inp_\iopp=(\Mmatrix^\mathsf{T} \cdot \vv, \mu'), \oracle_\iopp=f, \wit_\iopp=\pi_2(\bff))\in \LIN(\comfield,\cC,\pack),$$
    so that the witness has dimension $\codedim$.


    The error $\kappa_0^*$ is only present if $\degevalfield>0$. Otherwise, \cref{c:core_iop} (and thus \cref{c:virtual_iop} has one round less) and $\kappa_0^*$ is not part of the soundness errors. 
\end{theorem}
\begin{proof}
    \albert{Too vague?} The proof results from observing that \cref{c:virtual_iop} is the composition of an IOR $\Pi_1$ from $\LinBitsDual^{\mathsf{virtual}}(\bits, \pi, R,\cC)$ to $\LIN(\comfield,\cC, \pack)$, with an IOP $\Pi_2$ for $\LIN(\comfield,\cC, \pack)$.  In $\Pi_1$ the input is $(\inp=(\mu,\vv, \Mmatrix), \oracle = f, \wit=\bff).$ The output is $((A^{\mathsf{T}}\cdot \vv, \mu'), f, \pi_2(\bff))$. The IOR  $\Pi_1$ is exactly the same as \cref{c:core_iop}  with input $((\mu, \vv), f, \bfh)$, until Phase 4. Hence by the same arguments as in the proof of \cref{t:thm_core_IOPP}, it is perfectly complete and has errors $\kappa_0^*, (\kappa_{\red,i})_i$ with parameters set as $(R,\pi, \comfield,\codedim',\codedim_1,\codedim_2,\cC)$. The IOP $\Pi_2$ is perfectly complete and has round-by-round knowledge errors $\kappa_{\iopp, i}$ by assumption.
    %
    %\albert{Todo. Or is it too obvious to do it?}
\end{proof}


\begin{remark}[Instantiating \cref{c:virtual_iop}]
    \cref{c:virtual_iop} works in the exact same way as \cref{c:core_iop} until Phase 4, where an IOP for $\LIN(\comfield,\cC, \pack)$ is invoked on a slightly different instance-oracle-triple. The latter IOP is a standard combination of ring switching \cref{a:ring_switch_1} and WHIR \cref{whir}, plus perhaps additional processes to make verification time cheaper \cref{?}. The whole \cref{c:virtual_iop} can be instantiated as discussed in \cref{s:instantiation}, assuming this latter IOP is instantiated correctly (it is a standard IOP, so we do not discuss it further).  
\end{remark}





